\documentclass{../../note}

\usepackage{xcolor}
\usepackage{longtable}
\usepackage{array}
\usepackage{listings}

\lstset{
  basicstyle=\ttfamily\small,
  keywordstyle=\color{blue},
  commentstyle=\color{green!50!black},
  stringstyle=\color{red!60!black},
  numbers=left,
  numberstyle=\tiny,
  numbersep=6pt,
  breaklines=true,
  frame=single
}

\title{HUFE 数据结构第三讲}
\author{isomoes}

\begin{document}

\maketitle

\section{本次课定位}

本次课对应 5 次课安排中的第 3 次课，时长 3 小时。本讲在整体安排中原本还要兼顾查找与内排序，但为了让同学先把非线性结构中最综合的一章吃透，这一讲我们先\textbf{集中讲图}，把图的概念、存储、遍历、最小生成树和最短路径讲扎实。

图这一章的难点在于：它不像线性表那样只处理前后关系，也不像树那样有天然根结点，而是要同时处理\textbf{顶点、边、方向、权值、连通性和算法过程}。因此本讲的核心不是死记定义，而是建立一条稳定分析主线：\textbf{概念 - 表示 - 遍历 - 生成树 - 最短路径}。

\subsection{本次课目标}

\begin{itemize}
  \item 理解图、顶点、边、弧、度、路径、回路、连通图、强连通图等基础概念。
  \item 区分无向图、有向图、网、完全图、连通分量、强连通分量等常见术语。
  \item 掌握图的几种存储结构，尤其是邻接矩阵和邻接表，并了解十字链表与邻接多重表。
  \item 掌握深度优先搜索 DFS 与广度优先搜索 BFS 的思想、过程、复杂度与适用场景。
  \item 理解生成树、最小生成树的概念，掌握 Prim 和 Kruskal 算法的基本思想。
  \item 掌握最短路径问题的基本类型，重点理解 Dijkstra 算法的执行过程和使用条件。
  \item 能处理图这一章中常见的判断题、计算题、算法分析题和结构选择题。
\end{itemize}

\subsection{考试导向提醒}

图这一章在试卷中通常属于\textbf{高频 + 综合性较强}的内容，常见出题方式主要有五类：

\begin{itemize}
  \item \textbf{概念判断题}：如无向图与有向图、连通与强连通、生成树与最小生成树的区别。
  \item \textbf{存储结构题}：如邻接矩阵与邻接表的特点比较、度数统计、稀疏图与稠密图选型。
  \item \textbf{遍历过程题}：给定图和起点，写出 DFS 或 BFS 访问序列。
  \item \textbf{最小生成树题}：要求写出 Prim 或 Kruskal 的选边过程，或求最小生成树总权值。
  \item \textbf{最短路径题}：给定带权图，分析 Dijkstra 每一步的距离更新结果。
\end{itemize}

所以复习时不能只会背定义，一定要形成“\textbf{看结构、会表示、能遍历、能算路径、能比较算法}”的整体能力。

\subsection{3小时建议节奏}

\begin{center}
\begin{tabular}{|c|c|p{9cm}|}
\hline
阶段 & 时间 & 主要内容 \\
\hline
第1段 & 35分钟 & 图的基本概念：有向图、无向图、度、路径、连通、生成树等基础术语 \\
\hline
第2段 & 40分钟 & 图的存储结构：邻接矩阵、邻接表、十字链表、邻接多重表及特点比较 \\
\hline
第3段 & 45分钟 & 图的遍历：DFS、BFS 的思想、过程、复杂度与典型访问序列 \\
\hline
第4段 & 35分钟 & 生成树、最小生成树：Prim、Kruskal 的选边规则与适用图形 \\
\hline
第5段 & 25分钟 & 最短路径：Dijkstra 的核心思想、限制条件与典型题 \\
\hline
\end{tabular}
\end{center}

\section{图的基本概念}

\subsection{定义}

\textcolor{red}{图}是由顶点集 $V$ 和边集 $E$ 组成的二元组，记作：

\[
G = (V, E)
\]

其中，$V$ 是顶点的有限非空集合，$E$ 是连接顶点之间关系的边的有限集合。

和前面学过的结构相比，图最核心的特点是：\textbf{顶点之间的关系是多对多}。也正因为如此，图属于典型的非线性结构。

\subsection{有向图与无向图}

\begin{itemize}
  \item \textbf{无向图}：边没有方向，若顶点 $v_i$ 与 $v_j$ 相连，记作 $(v_i, v_j)$。
  \item \textbf{有向图}：边有方向，此时边也称为\textbf{弧}，记作 $\langle v_i, v_j \rangle$，其中 $v_i$ 是弧尾，$v_j$ 是弧头。
\end{itemize}

考试中要特别注意：

\begin{itemize}
  \item 在无向图中，边 $(v_i, v_j)$ 和 $(v_j, v_i)$ 表示同一条边。
  \item 在有向图中，弧 $\langle v_i, v_j \rangle$ 和 $\langle v_j, v_i \rangle$ 是两条不同的弧。
\end{itemize}

\subsection{简单图、完全图与网}

\begin{itemize}
  \item \textbf{简单图}：不存在重复边，也不存在顶点到自身的边。
  \item \textbf{完全图}：任意两个不同顶点之间都存在边。
  \item \textbf{网}：边或弧上带有权值的图，也叫带权图。
\end{itemize}

若无向完全图有 $n$ 个顶点，则边数为：

\[
\frac{n(n-1)}{2}
\]

若有向完全图有 $n$ 个顶点，则弧数为：

\[
n(n-1)
\]

这两个公式是非常高频的客观题考点。

\subsection{顶点的度}

\begin{itemize}
  \item 在\textbf{无向图}中，某顶点的\textbf{度}是与该顶点相关联的边数。
  \item 在\textbf{有向图}中，某顶点有两个量：\textbf{入度}和\textbf{出度}。
  \item 有向图中顶点的度满足：\textbf{度 = 入度 + 出度}。
\end{itemize}

两个常见结论：

\begin{enumerate}
  \item 在无向图中，所有顶点度数之和等于边数的 2 倍。
  \item 在有向图中，所有顶点入度之和等于所有顶点出度之和，且都等于弧数。
\end{enumerate}

\subsection{路径、回路与简单路径}

\begin{itemize}
  \item \textbf{路径}：从一个顶点到另一个顶点经过的顶点序列。
  \item \textbf{路径长度}：路径上边或弧的数目；若是带权图，也可能指路径上权值之和。
  \item \textbf{回路（环）}：起点和终点相同的路径。
  \item \textbf{简单路径}：除起点和终点可能相同外，其余顶点都不重复的路径。
  \item \textbf{简单回路}：除起点终点外，其余顶点不重复的回路。
\end{itemize}

\subsection{连通、连通图与强连通图}

\begin{itemize}
  \item 在无向图中，若两个顶点之间存在路径，则称这两个顶点是\textbf{连通}的。
  \item 若无向图中任意两个顶点都连通，则该图称为\textbf{连通图}。
  \item 连通图中的极大连通子图称为\textbf{连通分量}。
\end{itemize}

对有向图：

\begin{itemize}
  \item 若从顶点 $v_i$ 到 $v_j$ 和从 $v_j$ 到 $v_i$ 都存在路径，则称二者是\textbf{强连通}的。
  \item 若有向图中任意两个顶点都强连通，则称该图为\textbf{强连通图}。
  \item 有向图中的极大强连通子图称为\textbf{强连通分量}。
\end{itemize}

这里最容易混淆的是：\textbf{“连通”通常针对无向图，“强连通”通常针对有向图。}

\subsection{生成树与生成森林}

\begin{itemize}
  \item \textbf{生成树}：连通图的一个包含全部顶点的极小连通子图。
  \item \textbf{生成森林}：非连通图各个连通分量的生成树所组成的集合。
\end{itemize}

若一个连通无向图有 $n$ 个顶点，则它的任意一棵生成树都满足：

\begin{itemize}
  \item 恰好有 $n-1$ 条边；
  \item 去掉任意一条边后会变成非连通；
  \item 再增加任意一条原图中不在树内的边，就会形成回路。
\end{itemize}

\subsection{图的基础练习}

\begin{enumerate}
  \item 为什么图属于非线性结构？
  \item 无向图和有向图在边的表示上有什么区别？
  \item 无向完全图和有向完全图的边数公式分别是什么？
  \item 连通图和强连通图分别对应哪类图？
  \item 请判断：有向图中只要任意两个顶点之间存在单向路径，就称为强连通图。（\hspace{1cm}）
\end{enumerate}

\section{图的存储结构}

\subsection{为什么存储结构是重点}

图这一章中，很多算法都和存储结构直接相关。你选用不同的存储结构，会直接影响：

\begin{itemize}
  \item 判断两个顶点是否相邻是否方便；
  \item 找某个顶点全部邻接点是否高效；
  \item 空间开销大不大；
  \item 遍历和路径算法的时间复杂度如何计算。
\end{itemize}

因此图的学习一定不能只会背“有邻接矩阵和邻接表”，而要真正理解它们各自擅长什么。

\subsection{邻接矩阵}

\subsubsection{定义}

\textcolor{red}{邻接矩阵}是用一个二维数组来表示顶点之间邻接关系的存储方式。

若图有 $n$ 个顶点，则定义一个 $n \times n$ 的矩阵 $A$：

\begin{itemize}
  \item 对无向图，若 $v_i$ 和 $v_j$ 之间有边，则 $A[i][j] = 1$ 或对应权值，否则为 0 或 $\infty$。
  \item 对有向图，若从 $v_i$ 到 $v_j$ 有弧，则 $A[i][j] = 1$ 或对应权值，否则为 0 或 $\infty$。
\end{itemize}

\subsubsection{结构定义示意}

\begin{lstlisting}[language=C]
#define MaxVertexNum 100
typedef struct {
    char Vex[MaxVertexNum];
    int Edge[MaxVertexNum][MaxVertexNum];
    int vexnum, arcnum;
} MGraph;
\end{lstlisting}

\subsubsection{特点}

\textbf{优点：}
\begin{itemize}
  \item 判断两个顶点是否相邻非常快，可直接看矩阵元素，时间复杂度为 $O(1)$。
  \item 容易统计某顶点的度、入度、出度。
  \item 适合实现稠密图相关算法。
\end{itemize}

\textbf{缺点：}
\begin{itemize}
  \item 空间复杂度固定为 $O(|V|^2)$，对稀疏图不经济。
  \item 找某顶点全部邻接点时，通常要扫描整行或整列。
\end{itemize}

\subsubsection{高频判断}

\begin{enumerate}
  \item 无向图的邻接矩阵一定是对称矩阵。
  \item 无向图中第 $i$ 行元素之和等于顶点 $v_i$ 的度。
  \item 有向图中第 $i$ 行元素之和等于顶点 $v_i$ 的出度，第 $i$ 列元素之和等于其入度。
\end{enumerate}

\subsection{邻接表}

\subsubsection{定义}

\textcolor{red}{邻接表}是一种顺序存储与链式存储相结合的结构。它为每个顶点建立一个表头，再用链表记录该顶点的所有邻接点。

\subsubsection{结构定义示意}

\begin{lstlisting}[language=C]
typedef struct ArcNode {
    int adjvex;
    struct ArcNode *next;
} ArcNode;

typedef struct VNode {
    VertexType data;
    ArcNode *first;
} VNode, AdjList[MaxVertexNum];

typedef struct {
    AdjList vertices;
    int vexnum, arcnum;
} ALGraph;
\end{lstlisting}

\subsubsection{特点}

\textbf{优点：}
\begin{itemize}
  \item 空间复杂度约为 $O(|V| + |E|)$，更适合稀疏图。
  \item 找某顶点的全部邻接点方便。
  \item 在 DFS、BFS 中常更高效。
\end{itemize}

\textbf{缺点：}
\begin{itemize}
  \item 判断两个顶点是否直接相邻，不如邻接矩阵直观。
  \item 若要查找某一条边是否存在，通常需要遍历对应链表。
\end{itemize}

\subsubsection{无向图与有向图中的理解差异}

\begin{itemize}
  \item 在无向图的邻接表中，一条边通常会在两个顶点的链表中各出现一次。
  \item 在有向图的邻接表中，一条弧通常只会出现在弧尾对应的链表中一次。
\end{itemize}

因此：

\begin{itemize}
  \item 无向图中所有链表结点总数通常为 $2|E|$。
  \item 有向图中所有链表结点总数通常为 $|E|$。
\end{itemize}

\subsection{十字链表}

\subsubsection{为什么会有十字链表}

对于有向图，如果既希望方便找某顶点的\textbf{出边}，又希望方便找某顶点的\textbf{入边}，普通邻接表就不够直接了。

这时常用\textbf{十字链表}来存储有向图。

\subsubsection{核心思想}

十字链表中，每条弧结点同时挂到两条链上：

\begin{itemize}
  \item 一条是同一弧尾的出边链；
  \item 一条是同一弧头的入边链。
\end{itemize}

这样就可以同时方便地查找某顶点的出边和入边。

\subsubsection{适用对象}

\begin{itemize}
  \item 十字链表主要用于\textbf{有向图}。
  \item 它是“了解 + 识记”型考点，常考“适合存储什么图”“能方便求什么量”。
\end{itemize}

\subsection{邻接多重表}

\subsubsection{为什么会有邻接多重表}

在无向图中，一条边会分别出现在两个顶点的邻接表中，因此删除某条边时会涉及两个链表位置，不够自然。

为了解决这个问题，可以使用\textbf{邻接多重表}。

\subsubsection{核心思想}

邻接多重表让\textbf{每条边只对应一个边结点}，但这个边结点同时链接到与该边相关的两个顶点边表中。

它的好处是：

\begin{itemize}
  \item 对无向图中的边进行共享表示；
  \item 删除边、标记边等操作更方便。
\end{itemize}

\subsubsection{适用对象}

\begin{itemize}
  \item 邻接多重表主要用于\textbf{无向图}。
  \item 这也是常见的识记型考点。
\end{itemize}

\subsection{几种存储结构对比}

\begin{center}
\begin{tabular}{|>{\centering\arraybackslash}m{2.5cm}|>{\centering\arraybackslash}m{2.5cm}|>{\centering\arraybackslash}m{3.2cm}|>{\centering\arraybackslash}m{3.2cm}|}
\hline
存储结构 & 适合图类型 & 主要优点 & 主要缺点 \\
\hline
邻接矩阵 & 稠密图 & 判断相邻快，结构直观 & 空间开销大 \\
\hline
邻接表 & 稀疏图 & 节省空间，找邻接点方便 & 判断相邻不如矩阵快 \\
\hline
十字链表 & 有向图 & 同时便于找入边和出边 & 结构较复杂 \\
\hline
邻接多重表 & 无向图 & 一条边只存一次 & 结构较复杂 \\
\hline
\end{tabular}
\end{center}

\subsection{存储结构练习}

\begin{enumerate}
  \item 邻接矩阵为什么更适合稠密图？
  \item 邻接表为什么更适合稀疏图？
  \item 十字链表主要用于有向图还是无向图？
  \item 邻接多重表主要用于有向图还是无向图？
  \item 请判断：无向图的邻接矩阵一定是对称矩阵。（\hspace{1cm}）
\end{enumerate}

\section{图的遍历}

\subsection{为什么图的遍历比树更难}

树在遍历时天然没有回路，且通常有根结点；图则不同：

\begin{itemize}
  \item 图中可能存在回路；
  \item 图不一定连通；
  \item 同一个顶点可能通过多条路径被再次访问。
\end{itemize}

因此图的遍历必须解决一个关键问题：\textbf{防止重复访问}。这就是为什么 DFS 和 BFS 都需要设置访问标记数组 \texttt{visited}。

\subsection{深度优先搜索 DFS}

\subsubsection{思想}

\textcolor{red}{深度优先搜索}的基本思想是：从某个起始顶点出发，沿着一条路径尽可能往深处走，直到不能再继续，然后回退到上一个顶点，改走其他分支。

它的直观感觉类似于“走迷宫先一条路走到底，不通再回头”。

\subsubsection{递归写法示意}

\begin{lstlisting}[language=C]
void DFS(Graph G, int v) {
    visit(v);
    visited[v] = 1;
    for (w = FirstNeighbor(G, v); w >= 0; w = NextNeighbor(G, v, w)) {
        if (!visited[w]) {
            DFS(G, w);
        }
    }
}
\end{lstlisting}

\subsubsection{特点}

\begin{itemize}
  \item 常用递归实现，也可用栈实现。
  \item 在形式上与树的先序遍历比较接近。
  \item 适合分析连通性、寻找路径、构造 DFS 生成树等问题。
\end{itemize}

\subsubsection{复杂度}

\begin{itemize}
  \item 采用邻接矩阵时，时间复杂度通常为 $O(|V|^2)$。
  \item 采用邻接表时，时间复杂度通常为 $O(|V| + |E|)$。
\end{itemize}

\subsection{广度优先搜索 BFS}

\subsubsection{思想}

\textcolor{red}{广度优先搜索}的基本思想是：从起始顶点出发，先访问距离它最近的一层顶点，再访问更远的一层顶点，按“由近到远”逐层扩展。

它和树的层序遍历非常接近，因此通常借助\textbf{队列}实现。

\subsubsection{队列写法示意}

\begin{lstlisting}[language=C]
void BFS(Graph G, int v) {
    visit(v);
    visited[v] = 1;
    EnQueue(Q, v);
    while (!QueueEmpty(Q)) {
        DeQueue(Q, u);
        for (w = FirstNeighbor(G, u); w >= 0; w = NextNeighbor(G, u, w)) {
            if (!visited[w]) {
                visit(w);
                visited[w] = 1;
                EnQueue(Q, w);
            }
        }
    }
}
\end{lstlisting}

\subsubsection{特点}

\begin{itemize}
  \item 常借助队列实现。
  \item 与树的层序遍历联系最紧密。
  \item 在\textbf{无权图单源最短路径}问题中非常重要。
\end{itemize}

\subsubsection{复杂度}

\begin{itemize}
  \item 采用邻接矩阵时，时间复杂度通常为 $O(|V|^2)$。
  \item 采用邻接表时，时间复杂度通常为 $O(|V| + |E|)$。
\end{itemize}

\subsection{DFS 与 BFS 对比}

\begin{center}
\begin{tabular}{|c|c|c|}
\hline
比较项目 & DFS & BFS \\
\hline
核心思想 & 一条路走到底再回溯 & 一层一层向外扩展 \\
\hline
常用辅助结构 & 栈或递归 & 队列 \\
\hline
形式类比 & 树的先序遍历 & 树的层序遍历 \\
\hline
典型用途 & 连通性、路径搜索 & 分层访问、无权最短路 \\
\hline
\end{tabular}
\end{center}

\subsection{非连通图怎么遍历}

如果图不是连通图，只从一个顶点出发做一次 DFS 或 BFS，通常不能访问全部顶点。

因此完整遍历图时，通常要写成：

\begin{enumerate}
  \item 先把 \texttt{visited} 全部初始化为未访问；
  \item 依次检查每个顶点；
  \item 若该顶点尚未访问，就以它为起点再做一次 DFS 或 BFS。
\end{enumerate}

这一点很容易在选择题里出错。

\subsection{遍历练习}

\begin{enumerate}
  \item 图的遍历为什么必须设置访问标记？
  \item DFS 和 BFS 分别常借助什么辅助结构？
  \item 为什么说 BFS 和树的层序遍历很接近？
  \item 在邻接表上实现遍历时，时间复杂度通常是多少？
  \item 请判断：对任意图，从任一顶点出发做一次 BFS 一定能访问全部顶点。（\hspace{1cm}）
\end{enumerate}

\section{生成树与最小生成树}

\subsection{生成树再理解}

在连通无向图中，生成树的重点不是“任选一些边”，而是同时满足两个条件：

\begin{itemize}
  \item 覆盖全部顶点；
  \item 保持连通且不形成回路。
\end{itemize}

如果原图有 $n$ 个顶点，那么生成树必须有 $n-1$ 条边，这一点一定要记牢。

\subsection{最小生成树的定义}

\textcolor{red}{最小生成树（MST）}是连通加权无向图中，所有生成树里\textbf{边权值之和最小}的一棵。

这里有两个关键词：

\begin{itemize}
  \item 必须是\textbf{连通加权无向图}；
  \item 比较标准是\textbf{总权值最小}，不是边数最少，因为所有生成树边数本来都相同。
\end{itemize}

\subsection{Prim 算法}

\subsubsection{思想}

Prim 算法的思路可以概括为：

\begin{enumerate}
  \item 从某个顶点出发，先把它看成当前生成树的一部分；
  \item 每一步都从“已选顶点集合”和“未选顶点集合”之间，选一条权值最小的边；
  \item 把新顶点并入生成树；
  \item 重复直到所有顶点都加入。
\end{enumerate}

它的特点是：\textbf{逐点扩张，像一棵树慢慢长大。}

\subsubsection{过程示意}

\begin{lstlisting}[language=C]
void Prim(MGraph G, int start) {
    初始化已选顶点集合U = {start};
    while (U 中顶点数 < G.vexnum) {
        在一个端点属于U、另一个端点不属于U的所有边中
        选权值最小的一条;
        把新顶点和该边加入生成树;
    }
}
\end{lstlisting}

\subsubsection{适用场景}

\begin{itemize}
  \item Prim 更适合\textbf{稠密图}。
  \item 若采用邻接矩阵，写法较自然。
\end{itemize}

\subsection{Kruskal 算法}

\subsubsection{思想}

Kruskal 算法的思路可以概括为：

\begin{enumerate}
  \item 先把所有边按权值从小到大排序；
  \item 从最小边开始依次考察；
  \item 若当前边加入后不会形成回路，就把它选入；
  \item 直到选出 $n-1$ 条边为止。
\end{enumerate}

它的特点是：\textbf{逐边筛选，能选就选，但绝不成环。}

\subsubsection{过程示意}

\begin{lstlisting}[language=C]
void Kruskal(Graph G) {
    将所有边按权值从小到大排序;
    count = 0;
    for (每条边e按序考察) {
        if (加入e后不形成回路) {
            选入e;
            count++;
        }
        if (count == G.vexnum - 1)
            break;
    }
}
\end{lstlisting}

\subsubsection{适用场景}

\begin{itemize}
  \item Kruskal 更适合\textbf{稀疏图}。
  \item 它常与并查集思想联系起来理解，但考试中多数情况只要求掌握选边过程。
\end{itemize}

\subsection{Prim 与 Kruskal 对比}

\begin{center}
\begin{tabular}{|c|c|c|}
\hline
比较项目 & Prim & Kruskal \\
\hline
选边视角 & 从顶点集合向外扩张 & 从全局边集中依次选边 \\
\hline
核心限制 & 新边必须连接已选点和未选点 & 新边不能形成回路 \\
\hline
更适合 & 稠密图 & 稀疏图 \\
\hline
思维特征 & 长树 & 选边 \\
\hline
\end{tabular}
\end{center}

\subsection{最小生成树常见判断}

\begin{enumerate}
  \item 最小生成树一定是生成树，因此一定包含全部顶点。
  \item 最小生成树只针对连通加权无向图讨论。
  \item 最小生成树不一定唯一，但最小权值和相同。
  \item 生成树边数固定为 $n-1$，不是越少越好，而是在固定边数条件下总权值最小。
\end{enumerate}

\subsection{最小生成树练习}

\begin{enumerate}
  \item 生成树为什么一定没有回路？
  \item 连通无向图有 $n$ 个顶点时，其生成树有多少条边？
  \item Prim 和 Kruskal 的核心选边规则各是什么？
  \item 哪个算法更适合稀疏图？
  \item 请判断：最小生成树一定唯一。（\hspace{1cm}）
\end{enumerate}

\section{最短路径}

\subsection{最短路径问题的基本理解}

最短路径研究的是：从一个顶点到另一个顶点，怎样走总代价最小。

这里的“最短”可能有两种理解：

\begin{itemize}
  \item 在无权图中，通常指经过的边数最少；
  \item 在带权图中，通常指路径上边权值之和最小。
\end{itemize}

\subsection{BFS 与无权最短路径}

若图是无权图，或者所有边权都可视为相同，那么从某起点到其他顶点的最短路径问题，通常可由\textbf{BFS}解决。

原因是 BFS 按层次扩展，先到达的顶点一定是通过更少边数到达的。

这也是为什么很多选择题会把“BFS”和“无权图最短路径”放在一起考。

\subsection{Dijkstra 算法}

\subsubsection{适用条件}

\textcolor{red}{Dijkstra 算法}用于求解\textbf{单源最短路径}问题，即从一个源点到其他各顶点的最短路径。

但它有一个非常重要的限制：

\begin{center}
\textbf{Dijkstra 只适用于边权非负的图。}
\end{center}

若图中存在负权边，则不能直接使用 Dijkstra。

\subsubsection{核心思想}

它的基本思想可以概括为：

\begin{enumerate}
  \item 先确定源点到自身距离为 0，到其他点距离为初值；
  \item 每次从“尚未确定最短路径的顶点”中，选出当前距离最小的一个顶点；
  \item 把该顶点的最短距离正式确定下来；
  \item 再利用这个顶点去更新其邻接点的距离估计值；
  \item 重复直到所有顶点处理完毕。
\end{enumerate}

其中最关键的一句话是：\textbf{每一步都选当前暂时距离最小且未确定的顶点。}

\subsubsection{算法过程示意}

\begin{lstlisting}[language=C]
void Dijkstra(MGraph G, int v0) {
    初始化dist数组和visited数组;
    dist[v0] = 0;
    for (i = 0; i < G.vexnum; i++) {
        选择当前未确定且dist最小的顶点u;
        visited[u] = 1;
        for (u的每个邻接点v) {
            if (!visited[v] && dist[u] + weight(u, v) < dist[v]) {
                dist[v] = dist[u] + weight(u, v);
            }
        }
    }
}
\end{lstlisting}

\subsubsection{理解要点}

\begin{itemize}
  \item \texttt{dist[i]} 表示当前已知的源点到顶点 $i$ 的最短距离估计值。
  \item \texttt{visited[i]} 表示该顶点的最短距离是否已经最终确定。
  \item 一旦某顶点被选为当前最小距离顶点并确定下来，在非负权前提下，它的最短路径长度就不会再变。
\end{itemize}

\subsection{Floyd 算法简单了解}

如果题目问的是\textbf{任意两点之间}的最短路径，则更常见的算法是 Floyd 算法。

它属于多源最短路径算法，时间复杂度为 $O(|V|^3)$。在 HUFE 这一轮复习中，可以先把重点放在\textbf{Dijkstra 的思想和过程}上，而 Floyd 只作概念了解即可。

\subsection{最短路径练习}

\begin{enumerate}
  \item 无权图的单源最短路径通常可以用什么方法求？
  \item Dijkstra 解决的是单源最短路径还是多源最短路径？
  \item Dijkstra 为什么不能直接处理负权边？
  \item 若题目要求任意两点之间的最短路径，通常会想到哪种算法？
  \item 请判断：Dijkstra 算法适用于任意带权图。（\hspace{1cm}）
\end{enumerate}

\section{易混点总结}

\begin{enumerate}
  \item \textbf{无向图与有向图}：一个边无方向，一个弧有方向，表示方法不能混。
  \item \textbf{连通与强连通}：前者主要针对无向图，后者主要针对有向图。
  \item \textbf{邻接矩阵与邻接表}：前者判断相邻快，后者找邻接点更省空间。
  \item \textbf{DFS 与 BFS}：DFS 像走迷宫走到底，BFS 像一层一层扩散。
  \item \textbf{生成树与最小生成树}：最小生成树首先必须是生成树，然后才比较权值和。
  \item \textbf{Prim 与 Kruskal}：一个从点集出发扩张，一个从边集出发筛选。
  \item \textbf{BFS 与 Dijkstra}：前者常处理无权最短路，后者处理非负权单源最短路。
  \item \textbf{十字链表与邻接多重表}：十字链表偏有向图，邻接多重表偏无向图。
\end{enumerate}

\section{本讲知识框架}

建议把本讲内容整理成下面这条主线：

\begin{center}
\begin{tabular}{|>{\centering\arraybackslash}m{1.8cm}|>{\centering\arraybackslash}m{2.6cm}|>{\centering\arraybackslash}m{3.4cm}|>{\centering\arraybackslash}m{3.6cm}|}
\hline
主题 & 核心问题 & 典型内容 & 高频考点 \\
\hline
图基础 & 顶点如何建立关系 & 有向图、无向图、度、路径、连通 & 术语辨析、边数计算 \\
\hline
存储结构 & 图在内存中怎样表示 & 邻接矩阵、邻接表 & 稠密/稀疏图选型、度统计 \\
\hline
图遍历 & 如何不重不漏访问图 & DFS、BFS & 访问序列、复杂度、辅助结构 \\
\hline
生成树 & 如何保持连通又无环 & 生成树、MST & 边数、Prim、Kruskal \\
\hline
最短路径 & 怎样走代价最小 & BFS、Dijkstra、Floyd & 无权最短路、负权限制 \\
\hline
\end{tabular}
\end{center}

\section{典型例题}

\subsection{例题1：无向完全图边数}

\textbf{题目：}一个无向完全图有 6 个顶点，则共有多少条边？

\textbf{分析：}无向完全图边数公式为 $\frac{n(n-1)}{2}$。

\textbf{解：}
\[
\frac{6 \times 5}{2} = 15
\]

\textbf{答案：}15 条边。

\subsection{例题2：邻接矩阵中的度}

\textbf{题目：}某无向图的邻接矩阵中，第 3 行元素之和为 4，则顶点 $v_3$ 的度是多少？

\textbf{分析：}无向图邻接矩阵中某一行元素之和就是该顶点的度。

\textbf{答案：}4。

\subsection{例题3：遍历辅助结构}

\textbf{题目：}图的广度优先搜索通常借助什么线性结构实现？

\textbf{分析：}广度优先要求按先到先处理的顺序逐层扩展，符合先进先出。

\textbf{答案：}队列。

\subsection{例题4：生成树边数}

\textbf{题目：}一个连通无向图有 8 个顶点，则其任意生成树有多少条边？

\textbf{分析：}连通无向图的生成树边数恒为 $n-1$。

\textbf{答案：}7 条边。

\subsection{例题5：Prim 与 Kruskal 的区别}

\textbf{题目：}Prim 和 Kruskal 都用于求最小生成树，它们的主要思路有什么不同？

\textbf{分析：}Prim 从顶点集合向外扩张，Kruskal 从全局边集按权值筛选且避免成环。

\textbf{答案：}Prim 是“逐点扩张”，Kruskal 是“逐边选择”。

\subsection{例题6：Dijkstra 的使用条件}

\textbf{题目：}某带权图中存在负权边，是否还能直接使用 Dijkstra 算法？

\textbf{分析：}Dijkstra 建立在“已确定的最短距离不再变化”的非负权前提上。

\textbf{答案：}不能。

\subsection{例题7：BFS 解决什么最短路}

\textbf{题目：}无权图中，从某一源点到其余顶点的最短路径问题，通常优先使用 DFS 还是 BFS？

\textbf{分析：}无权图最短路径等价于边数最少，而 BFS 按层次扩展。

\textbf{答案：}BFS。

\section{课堂练习}

\subsection{基础练习}

\begin{enumerate}
  \item 图和树相比，为什么说图的关系更复杂？
  \item 有向图中顶点的度、入度、出度三者有什么关系？
  \item 为什么无向图邻接矩阵一定是对称矩阵？
  \item 邻接表为什么通常更适合稀疏图？
  \item DFS 和 BFS 在访问策略上的根本差异是什么？
\end{enumerate}

\subsection{判断与选择}

\begin{enumerate}
  \item 图中任意两个顶点之间都存在边，则该图一定是完全图。（\hspace{1cm}）
  \item 无向图中所有顶点度数之和等于边数。（\hspace{1cm}）
  \item 有向图中所有顶点入度之和等于所有顶点出度之和。（\hspace{1cm}）
  \item BFS 通常借助栈实现。（\hspace{1cm}）
  \item 最小生成树一定包含原图全部顶点。（\hspace{1cm}）
  \item Dijkstra 可以直接处理负权边。（\hspace{1cm}）
\end{enumerate}

\subsection{计算与分析}

\begin{enumerate}
  \item 一个无向完全图有 5 个顶点，求边数。
  \item 一个有向完全图有 4 个顶点，求弧数。
  \item 一个连通无向图有 10 个顶点，求任意生成树的边数。
  \item 为什么在图遍历中必须设置 \texttt{visited} 数组？
  \item Prim 算法和 Kruskal 算法都求最小生成树，但各自的选边标准有什么不同？
  \item 为什么说 BFS 常可以求无权图的单源最短路径？
\end{enumerate}

\subsection{综合小题}

\textbf{题目：}某校园系统要处理三个任务：

\begin{itemize}
  \item 任务 A：表示校园楼宇之间的道路连接关系，并经常判断两栋楼是否直接相连；
  \item 任务 B：从某一教学楼开始，按“距离越来越远”的层次依次检查周边楼宇；
  \item 任务 C：在所有楼宇都连通的前提下，选择一组总造价最低的道路方案。
\end{itemize}

请分别指出任务 A、任务 B、任务 C 对应的核心数据结构或算法思想，并说明理由。

\section{课后小测参考答案}

\begin{enumerate}
  \item 因为图中一个顶点可以和多个顶点发生关系，整体是多对多结构。
  \item 有向图中顶点的度 = 入度 + 出度。
  \item 因为无向图中边 $(v_i, v_j)$ 与 $(v_j, v_i)$ 表示同一条边，所以矩阵关于主对角线对称。
  \item 因为空间复杂度更接近 $O(|V| + |E|)$，对稀疏图更经济。
  \item DFS 强调先深入再回溯，BFS 强调按层次逐步扩展。
  \item 第1题对；第2题错；第3题对；第4题错；第5题对；第6题错。
  \item 无向完全图边数为 $\frac{5 \times 4}{2} = 10$。
  \item 有向完全图弧数为 $4 \times 3 = 12$。
  \item 生成树边数为 $10 - 1 = 9$。
  \item 因为图中可能存在回路和重复到达路径，若不标记已访问顶点就可能重复访问甚至死循环。
  \item Prim 每次选连接已选点集和未选点集的最小边；Kruskal 每次从全局最小边开始选且不允许成环。
  \item 因为 BFS 按边数从少到多逐层扩展，先访问到的顶点对应的路径边数最少。
\end{enumerate}

综合小题参考：任务 A 若经常判断两栋楼是否直接相连，可优先考虑邻接矩阵；任务 B 对应 BFS，因为是按层次向外扩展；任务 C 对应最小生成树思想，可用 Prim 或 Kruskal 算法。

\subsection{分节练习参考答案}

\textbf{图的基础练习参考：}

\begin{enumerate}
  \item 因为顶点之间是多对多关系，不是一对一或一对多关系。
  \item 无向图用无序对表示边，有向图用有序对表示弧。
  \item 无向完全图为 $\frac{n(n-1)}{2}$，有向完全图为 $n(n-1)$。
  \item 连通图对应无向图，强连通图对应有向图。
  \item 错。强连通要求任意两点互相可达，而不是只存在单向路径。
\end{enumerate}

\textbf{存储结构练习参考：}

\begin{enumerate}
  \item 因为邻接矩阵空间固定为 $O(|V|^2)$，稠密图中空间浪费相对较少，且判断相邻很快。
  \item 因为邻接表空间更接近 $O(|V| + |E|)$，当边少时更节省空间。
  \item 十字链表主要用于有向图。
  \item 邻接多重表主要用于无向图。
  \item 对。
\end{enumerate}

\textbf{遍历练习参考：}

\begin{enumerate}
  \item 因为图中可能有回路，也可能通过不同路径多次到达同一顶点。
  \item DFS 常借助栈或递归，BFS 常借助队列。
  \item 因为它们都按距离起点越来越远的层次顺序访问结点或顶点。
  \item 在邻接表上通常为 $O(|V| + |E|)$。
  \item 错。若图不连通，一次 BFS 不能访问全部顶点。
\end{enumerate}

\textbf{最小生成树练习参考：}

\begin{enumerate}
  \item 因为生成树要求在保持连通的同时不形成回路。
  \item $n-1$ 条边。
  \item Prim 是从已选点集向外选最小边，Kruskal 是按权值从小到大选边且避免成环。
  \item Kruskal 更适合稀疏图。
  \item 错。若最小权值组合不唯一，则最小生成树也可能不唯一。
\end{enumerate}

\textbf{最短路径练习参考：}

\begin{enumerate}
  \item BFS。
  \item 单源最短路径。
  \item 因为负权边会破坏“当前最小距离一经确定就不再变化”的基础。
  \item Floyd 算法。
  \item 错。Dijkstra 只适用于边权非负的图。
\end{enumerate}

\section{本讲小结}

本讲真正要建立的不是一堆零散名词，而是一种看图问题的稳定方法：\textbf{先看顶点和边的关系，再看用什么结构存，再看如何遍历，最后再上升到生成树和最短路径问题}。

如果说前两讲分别在建立线性结构和树形结构的理解，那么图这一讲就是把“多对多关系 + 算法过程”真正引进来。这里最需要反复练习的四个点分别是：\textbf{邻接矩阵与邻接表的区别、DFS 和 BFS 的访问思想、Prim 与 Kruskal 的选边逻辑、Dijkstra 的使用条件与更新过程}。

下一步如果继续补 class3 后半部分内容，就可以在此基础上接上\textbf{查找和内排序}，形成第 3 次课的完整数据结构收尾讲义。

\end{document}
